\documentclass[11pt]{article}
\usepackage[a4paper,margin=25mm]{geometry}
\usepackage{amsmath,amssymb,amsthm,booktabs,xurl,hyperref}
\hypersetup{hidelinks}
\newtheorem{theorem}{Theorem}
\newtheorem{lemma}[theorem]{Lemma}
\newtheorem{proposition}[theorem]{Proposition}
\newcommand{\Q}{\mathbb Q}
\title{One polynomial denominator for a progression of contact equations}
\author{Z diophantine in Q project}
\date{24 November 2026}
\begin{document}
\maketitle
\begin{abstract}
Polynomial division replaces each contact combination register by a
shared coefficient representation of its denominator. Applied to the
retained contact progression and its sixteenth endpoint, this gives
complete systems with 963 variables and 818 quadratic equations in the
general route, and 1140 variables and 991 equations in the specified
route. These are exact rational changes of coordinates from the v102
systems: they preserve their entire rational zero sets, including all
exceptional inverse-weight values. All thirteen or eleven integer
predicates and every arithmetic premise remain unchanged.
\end{abstract}

\section{The elementary reduction}

Fix rational constants $A,B,C,D,E,F$ with $C\ne0$. Let $\kappa$
be a nonzero rational coordinate, or the constant one. Suppose existing
coordinates $u,f$ satisfy
\begin{equation}\label{square}
                            u^2=\kappa f.
\end{equation}
The two contact equations, with private coordinates $W,t$, are
\begin{align}
 \mathcal C&=Auf+(Bu+Cf)(W+\kappa)-\kappa t=0,\label{comb}\\
 \mathcal R&=(Du+Ef+F\kappa)(W+\kappa)+ft=0.\label{rel}
\end{align}
Assume that $t$ occurs nowhere else in the complete system. Occurrences
of $W$ elsewhere are permitted, and will be changed affinely below.
Define the quadratic expressions
\begin{align}
 P_*&=Cf^2+Buf+E\kappa f+D\kappa u+F\kappa^2,\label{P}\\
 R_*&=A(B^2-CE)uf+A(BE-CD)\kappa f\nonumber\\
    &\hspace{13mm}+A(BD-CF)\kappa u+ABF\kappa^2.\label{R}
\end{align}
Introduce the affine replacement
\begin{equation}\label{V}
                  V=C^2(W+\kappa)+ACu-AB\kappa.
\end{equation}
If $p$ is an already represented coordinate satisfying
\begin{equation}\label{p}
                             \kappa p=P_*,
\end{equation}
then replace the two contact equations by the single quadratic equation
\begin{equation}\label{new}
                                  pV+R_*=0.
\end{equation}
Here $V$ replaces $W$, rather than being added as another coordinate.

\begin{lemma}\label{lemma}
Under \eqref{square} and \eqref{p}, the replacement preserves the full
rational solution set bijectively after deleting $t$. Neither $V$ nor
$W$, $p$, or $P_*$ is required to be nonzero.
\end{lemma}
\begin{proof}
There is the exact polynomial identity
\begin{equation}\label{identity}
 P_*V+\kappa R_*-C^2(\kappa\mathcal R+f\mathcal C)
       =AC(Bf+D\kappa)(u^2-\kappa f),
\end{equation}
where $V$ is given by \eqref{V}. The verifier expands both sides in
independent symbols and checks this identity over $\mathbb Z$.
An old solution therefore gives \eqref{new}, since $\kappa\ne0$.
Conversely, recover the old coordinates uniquely by
\begin{equation}\label{recovery}
 W=\frac{V-ACu-(C^2-AB)\kappa}{C^2},\qquad
 t=\frac{Auf+(Bu+Cf)(W+\kappa)}{\kappa}.
\end{equation}
The combination equation holds by definition; \eqref{identity} now
gives $C^2\kappa\mathcal R=0$, hence the other old equation. These
maps divide only by the stated nonzero constant $C$ and the already
nonzero $\kappa$. Every occurrence of $W$ in another quadratic row is
changed by the affine first expression in \eqref{recovery}, so its
degree remains at most two.
\end{proof}

For motivation, in the unscaled coordinate $x=u/\kappa$ the relevant
denominator is
\[
 P(x)=Cx^4+Bx^3+Ex^2+Dx+F.
\]
Polynomial division of $Ax^5$ by $P(x)$ has quotient
$Ax/C-AB/C^2$. Adding this quotient to $1+W/\kappa$ and
multiplying by $\kappa C^2$ gives \eqref{V}. The remainder is only cubic, so
\eqref{square} represents it quadratically through $uf,f,u,1$.
The division is an identity of polynomials with a fixed nonzero leading
coefficient; it does not require that $P(x)$ be nonzero.

\section{Sharing the denominator coefficients}

Suppose a finite set $J\subset\Q$ labels contacts with retained affine
coordinate expressions
\begin{equation}\label{progression}
 u(j)=u_0+2jv+j^2w,\qquad
 f(j)=f_0+ja_1+j^2a_2+j^3a_3+j^4a_4.
\end{equation}
The existing system ensures $u(j)^2=\kappa f(j)$. In particular this
follows from $u_0^2=\kappa f_0$ and the four retained equations
\[
 \kappa a_1=4u_0v,\quad
 \kappa a_2=2u_0w+4v^2,\quad
 \kappa a_3=4vw,\quad \kappa a_4=w^2.
\]
Indeed the discrepancy in the identity for general $j$ is the sum of
the root discrepancy and the four discrepancies weighted by
$-j,-j^2,-j^3,-j^4$.

Form the degree-eight polynomial in an indeterminate $T$,
\begin{equation}\label{common}
 \Pi(T)=C f(T)^2+B u(T)f(T)+E\kappa f(T)
                          +D\kappa u(T)+F\kappa^2.
\end{equation}
Add nine common coordinates $p_0,\ldots,p_8$ and the nine quadratic rows
\begin{equation}\label{coeff}
                      \kappa p_j=[T^j]\Pi(T),\qquad 0\le j\le8.
\end{equation}
All coefficients of $u(T),f(T)$ are existing coordinates or affine
expressions in those coordinates, so every right side is quadratic.
For a contact at index $j$, use $p(j)=\sum_{h=0}^8p_hj^h$ in
\eqref{new}. This is an affine expression in the new coordinates and
therefore preserves degree two.

\begin{theorem}\label{general}
For $r$ contacts of the form \eqref{comb}--\eqref{rel} along
\eqref{progression}, with $r$ private combination registers occurring
only in their own two equations, the construction deletes those $r$
registers and replaces them by nine common registers. It saves $r-9$
variables and $r-9$ quadratic rows. The rational zero sets before and
after the construction are bijective under the explicit maps above.
The index set may contain zero or negative indices; the spacing can
vanish; and all inverse weights, including zero values, are covered.
\end{theorem}
\begin{proof}
Given an old solution, \eqref{coeff} has a unique solution because
$\kappa\ne0$. Apply Lemma~\ref{lemma} to each contact. Conversely,
those common equations imply \eqref{p} at every contact index, so the
lemma reconstructs all deleted coordinates and rows. Each forward and
backward coordinate map is rational and inverse to the other. The
change from $2r$ rows to $r+9$ rows gives the stated count.
\end{proof}

\section{Application to the complete v102 systems}

The v102 endpoint theorem \cite{endpoint} already identifies the second
contact of the root block as index 16 of the common progression. Its
square and fourth-power registers have been recharted to the same
gauge as all the first contacts. Set $u_0=\mathtt{int\_b\_60}$,
$f_0=\mathtt{int\_b\_61}$,
$u_{16}=\mathtt{int\_b\_69}$,
$f_{16}=\mathtt{int\_b\_70}$, and use
\[
 v=\frac{u_{16}-u_0-256w}{32},\qquad
 a_1=\frac{f_{16}-f_0-256a_2-4096a_3-65536a_4}{16}.
\]
These affine expressions are exactly those already retained by v102;
no new choice of contacts is made.

For the general route, $\kappa=1$ and
\begin{align*}
 A&=3813373161796,&B&=122376041313918,\\
 C&=-35893742479489,&D&=141822723674536,\\
 E&=-192222345736537,&F&=-39896049934224.
\end{align*}
Its first contacts have indices
$0,1,2,3,4,5,6,7,8,10,11,12,13$; adding the root second contact
at index 16 gives $r=14$. For the specified route,
$\kappa=\rho=\mathtt{int\_b\_162}$ is already guarded by
$\rho\mathtt{int\_b\_163}=1$, and
\begin{align*}
 A&=7878870169,&B&=1229741066709366,\\
 C&=-6786192246520,&D&=476802068222265625,\\
 E&=-42042541279847500,&F&=-1738195040039062500.
\end{align*}
Its first contact indices are
$0,1,2,3,4,5,6,7,8,10,11$; the same endpoint adds the twelfth
contact. Every displayed constant is checked against the literal
source rows, including the row normalization factors.

The deleted coordinates are the first-contact combination registers
$\mathtt{int\_block\_62}$ and the root second-contact register
$\mathtt{int\_b\_71}$. The inverse registers retain their names but
now store $V$ in \eqref{V}; this changed interpretation is recorded
explicitly in the complete tables. Other contact lists are unchanged.
An exhaustive occurrence check verifies that each deleted coordinate
occurs in exactly its combination and relation rows. Every changed
occurrence of each inverse register is included, in particular its
inverse-weight and gate rows.

\begin{center}
\begin{tabular}{lrrrr}
\toprule
Route & Old variables & New variables & Old rows & New rows\\
\midrule
General & 968 & 963 & 823 & 818\\
Specified & 1143 & 1140 & 994 & 991\\
\bottomrule
\end{tabular}
\end{center}
The ordinary-$abc$ construction retains 695 variables and 590 quadratic
rows. The sums of squared residuals
are polynomials of degree exactly four: a retained outer row has a
nonzero quadratic part, and a sum of real squares of homogeneous
quadratic parts cannot cancel identically.

The exact changes preserve all thirteen integer predicates of the
general route and all eleven of the specified route. They change no
coefficient parameter or arithmetic interface. The general route's
sufficient parameter remains ineffective. All conclusions about
undecidability are still under the retained arithmetic premises of the
integer interfaces and the externally cited Jones coding theorem.
This paper establishes no unconditional undecidability conclusion,
no global minimality result, and no quartic bound close to 263.

\section{Reproduction and verification scope}

The deterministic compiler and independent checker are
\begin{quote}\small
\nolinkurl{undecidability/quadratic_reduction_v103/build_contact_denominator.py}\\
\nolinkurl{undecidability/quadratic_reduction_v103/verify_contact_denominator.py}.
\end{quote}
The complete row tables and adjacent certificates are
\begin{quote}\small
\nolinkurl{denominator_general_quadrics.json},\quad
\nolinkurl{denominator_general_certificate.json},\\
\nolinkurl{denominator_specified_quadrics.json},\quad
\nolinkurl{denominator_specified_certificate.json}.
\end{quote}
The checker imports no compiler. It proves \eqref{identity} symbolically,
pins the retained root and four progression rows, independently selects
the contact indices and constants, checks every omitted-coordinate
occurrence, derives each changed row from the recovery map, and checks
every unchanged row literally. It verifies the nine coefficient equations,
complete coordinate support, coefficient domain, degree, counts, parameter
and integer-predicate preservation, and input/output SHA-256 hashes.
Its result is a finite algebraic audit and does not independently prove
the inherited arithmetic interfaces.

\begin{thebibliography}{9}
\bibitem{endpoint} Z diophantine in Q project,
\emph{Contact endpoint reuse}, 23 November 2026,
\nolinkurl{undecidability/contact_endpoint_v102/contact_endpoint_reuse.tex}.
\bibitem{progression} Z diophantine in Q project,
\emph{Affine contact progressions}, 23 November 2026,
\nolinkurl{undecidability/general_contact_v94/affine_contact_progressions.tex}.
\end{thebibliography}
\end{document}
